%%%%%%%%%%%%%%%%%%%%%%%%%%%%%%%%%%%%%%%%%
% Twenty Seconds Resume/CV
% LaTeX Template
% Version 1.0 (14/7/16)
%
% Original author:
% Carmine Spagnuolo (cspagnuolo@unisa.it) with major modifications by
% Vel (vel@LaTeXTemplates.com) and Harsh (harsh.gadgil@gmail.com)
%
% License:
% The MIT License (see included LICENSE file)
%
%%%%%%%%%%%%%%%%%%%%%%%%%%%%%%%%%%%%%%%%%

%----------------------------------------------------------------------------------------
%	PACKAGES AND OTHER DOCUMENT CONFIGURATIONS
%----------------------------------------------------------------------------------------

\documentclass[letterpaper]{twentysecondcv} % a4paper for A4

% Command for printing skill overview bubbles
\newcommand\skills{
	\smartdiagram[bubble diagram]{
        \textbf{Ingeniería}\\\textbf{en}\\\textbf{Software},
        \textbf{Programación},
        \textbf{~~~~~DevOps~~~~~~~~},
        \textbf{~~~~~Ciencia~~~~~~~}\\\textbf{~~de Datos~~},
        \textbf{~Arquitectura~}
    }
}

% Technical skill bars
\programming{{Terraform $\textbullet$ Crossplane $\textbullet$ FluxCD / 5},{AzDO $\textbullet$ GitLab CI$\slash$CD $\textbullet$ GitHub Actions / 5},{C $\textbullet$ Python $\textbullet$ JS$\slash$TS $\textbullet$ Java $\textbullet$ \LaTeX / 6},{AWS: ec2 $\textbullet$ s3 $\textbullet$ rds $\textbullet$ eks$\ldots$ / 5}}

% Soft skills bullet points
\softskills{
    Comunicativo $\textbullet$ Persistente $\textbullet$ Directo\\
    Creativo $\textbullet$ Profesional $\textbullet$ Detallista\\
    Puntual $\textbullet$ Estudioso $\textbullet$ Trabajo en Equipo
}

% Education
\education{
\textbf{Ingeniería Física} \\
Universidad de Santiago de Chile \\
2017 - 2022 | Santiago, Chile
}

%----------------------------------------------------------------------------------------
%	 PERSONAL INFORMATION
%----------------------------------------------------------------------------------------
% If you don't need one or more of the below, just remove the content leaving the command, e.g. \cvnumberphone{}

\cvname{Andoni\\Vizcaya Rodríguez \\ \color{rpm_text}\footnotesize{19.287.552-8}} % Your name
\cvjobtitle{ Físico$\slash$Ingeniero en Software } % Job
% title/career

\cvlinkedin{Andoni Vizcaya}
\cvgithub{andonivizcaya} % It was supposed to be my PornHub but is ok :(
\cvnumberphone{+56 9 4411 9108} % Phone number
\cvsite{08/10/1995} % More wia
\cvmail{andonivizcaya@gmail.com} % Email address

%----------------------------------------------------------------------------------------

\begin{document}

\makeprofile % Print the sidebar

%----------------------------------------------------------------------------------------
%	 EXPERIENCE
%----------------------------------------------------------------------------------------

\section{Experiencia}

\begin{twenty} % Environment for a list with descriptions
\twentyitem
    	{\small May 2023 -}
		{\small Presente}
        {Senior DevOps Engineer}
        {\href{https://sovos.com}{Sovos}}
        {\small Las Condes, Chile}
        {
        \begin{itemize}
            \item IaC en Terraform y Crossplane para despliegue de infraestructura
            \item Pipelines CI/CD en AzDO
            \item Clusters k8s multi-tenat usando herramientas modernas
            \item Arquitectura de software para aplicaciones nativas en cloud
            \item Asumir la responsabilidad de \textit{Platform Engineer} dentro del equipo
            \item Migración y rediseño de servicios y aplicaciones on-premise a la nube
            \item Infraestructura y desarrollo para aplicaciones ETL
        \end{itemize}}
        \\

\twentyitem
    	{\small Dic 2021 -}
		{\small May 2023}
        {Desarrollador de Software}
        {\href{https://pleyasoft.com/}{PleyaSoft}/\href{https://cencosud.com}{Cencosud}}
        {\small Ñuñoa, Chile}
        {
        \begin{itemize}
            \item Aplicación full-stack de autoservicio Cencosud (Node/React/Native)
            \item Meta-aplicación para realizar tests unitarios y end to end
            \item Dashboards usando herramientas de data science en python
            \item Responsable de despliegues de la aplicación en AWS
        \end{itemize}}
        \\

\twentyitem
    	{\small Jun 2021 -}
		{\small Dic 2021}
        {Desarrollador de Software}
        {\href{https://athos.cl/}{Athos}/\href{https://editrade.cl}{Editrade}}
        {\small Providencia, Chile}
        {
        \begin{itemize}
            \item Desarrollo de módulos para Odoo en python y C
            \item Aplicación de monitoreo y automatización de respaldos de DBs
            \item Hardware usando arduino, raspberry pi y más
        \end{itemize}}
        \\

\twentyitem
    	{\small Mar 2019 -}
		{\small Mar 2020}
        {Asistente de investigación}
        {\href{https://uchile.cl/}{Universidad de Chile}/\href{https://dfi.uchile.cl/}{dfi (FCFM)}}
        {\small Santiago, Chile}
        {
        \begin{itemize}
            \item Investigación en Teoría de Turbulencia de Ondas (WT)
            \item Confección de setup experimental, toma y análisis de datos
            \item Utilización de ML para crear modelos predictivos (C++/python)
        \end{itemize}}

	%\twentyitem{<dates>}{<title>}{<location>}{<description>}
\end{twenty}

%----------------------------------------------------------------------------------------
%	 RESEARCH
%----------------------------------------------------------------------------------------
\section{Proyectos Técnicos}
\begin{twenty}
	\twentyitem
    	{\small Abr 2024 -}
		{\small Presente}
        {Hydrogen Orbital}
        {\href{https://github.com/andonivizcaya/hydrogen_orbital}{hydrogen\_orbital}}
        {\small Santiago}
        {
        {\begin{itemize}
            \item Visualización de orbitales atómicos usando raylib con C puro
            \item Proyecto compilado a Linux, MacOS, Windows y WebAssembly
            \item Wasm corriendo en \href{https://andonivizcaya.github.io/hydrogen_orbital/}{\textbf{hydrogen\_orbital}} con CI en github actions
            \item Trabajando en implementación de ágebra geométrica para rotores
		\end{itemize}
        }}
\end{twenty}

\section{Idiomas}
\begin{itemize}
    \item \textbf{Español}: nativo
    \item \textbf{Inglés}: avanzado
\end{itemize}


\section{Hobbies}
\begin{itemize}
    \item \textbf{Deporte} (trekking/hiking, natación, escalada)
    \item \textbf{Música} (composición y como usuario)
    \item \textbf{Ciencia} (física > matemáticas > electrónica > vim)
\end{itemize}

\end{document}

